\section{Weighted disjointness in bounded space}\label{sec:separator}

The parent schedule produces two arrays of records, each consisting of
an actual weight and a top-mixer mask. We want one record from each
array whose weights sum to a given target and whose masks are disjoint. A separating set
turns this into an ordinary weight search: retain the left records whose
sets lie inside it, and the right records whose sets lie outside it.
Every pair of retained records is then disjoint. Sorting their weights
finds an exact complementary pair, if one exists.

The simplest illustration uses a block of $g$ indices, with $g$ divisible
by four. Suppose each side of a fixed disjoint pair occupies $g/4$
indices. A uniformly sampled set of size $g/2$ separates this pair with
probability
\[
 \frac{\binom{g/2}{g/4}}{\binom{g}{g/2}}.
\]
Indeed, the sampled set must include the left $g/4$ indices, exclude
the right $g/4$ indices, and choose its remaining $g/4$ elements from
the $g/2$ unused indices. Sampling a family whose size is a constant
times the reciprocal of this probability gives a constant chance of
covering the fixed pair. An individual record can occur in several
members of this family; those occurrences determine the work.

We will apply this idea on four blocks. The left and right cardinalities
may differ, so the sampled set size is biased to control occurrences on
both sides. We first specify the distribution. We then show how to
search its useful index tuples without storing all tuples or testing
all pairs of records.

This construction follows the separating-set approach of Nederlof and
W\k{e}grzycki~\cite[full version, Section~6]{NW21}, which builds on the
separating collections of Fomin, Lokshtanov, Panolan, and
Saurabh~\cite[Section~4.2, Definition~4.3 and Lemma~4.5]{FLPS16}.
We use the affine bias from Nederlof and W\k{e}grzycki's
construction~\cite[full version, Lemma~6.6]{NW21}, with the
cardinality-range extension of its entropy estimate by Belova, Chukhin,
Kulikov, and Mihajlin~\cite[full version, Section~4.7]{BCKM24}.
The rounding, traversal, and exact estimate required here are proved below.

\subsection{Inputs and sampled families}

Recall the parameters $u,m,g,C$ from \eqref{eq:parameters}. The four
blocks in the representation layout give the partition
\[
 M_0=T_1\sqcup T_2\sqcup T_3\sqcup T_4,\qquad |T_i|=g.
\]
Write $k=k_0/m\in[1/6,1/4]$, where $k_0$ is the prescribed left
mask size. To recall the balanced block counts explicitly, write
$k_0=4d+r$ with $0\le r<4$ and put $c_i=d+1$ for $i\le r$ and
$c_i=d$ otherwise. Then $\sum_i c_i=k_0$ and $|c_i-kg|<1$.
The left array $\mathcal L$ and right array
$\mathcal R$ are nonempty and have lengths at most $C$. Their records
are $(x,D)$ and $(y,E)$, respectively, with $D,E\subseteq M_0$ and
\begin{equation}\label{sep:cardinalities}
 |D\cap T_i|=c_i,\qquad |E\cap T_i|=g/2-c_i
 \quad(1\le i\le4).
\end{equation}
These conditions do not assert that an arbitrary pair is disjoint.
For the transformed target $t'$, the required output is YES only if some
pair satisfies
\begin{equation}\label{sep:target}
 D\cap E=\varnothing,\qquad x+y=t'.
\end{equation}
Array lengths count occurrences; distinct records are not required.

The actual weights and target have magnitude below
$2^{b+\lceil\log(n+2)\rceil+4}$. The exact arithmetic, masks, and
array indices used here fit in a constant number of the words specified
in Section~\ref{sec:introduction}; all stored tables and retained
random choices count as workspace. The space bound for this check
will hold on every outcome of the randomness.

Let
\[
 \alpha=\log 3,\qquad
 \alpha_m=\frac{\lfloor\log(3^m)\rfloor}{m},\qquad
 e=\left\lfloor g\left(\frac12+
              \alpha_m\left(k-\frac14\right)\right)\right\rfloor.
\]
The logarithm floor in $\alpha_m$ is the bit length of the integer
$3^m$ minus one; no real-number oracle is needed. For each block set
\begin{equation}\label{sep:families}
 \nu_i=\left\lceil
  10\binom{g}{e}\bigg/\binom{g/2}{e-c_i}
 \right\rceil,
 \qquad \mathcal F_i=(S_{i,1},\ldots,S_{i,\nu_i}).
\end{equation}
All $S_{i,j}$ are independent, and $S_{i,j}$ is a uniform $e$-element
subset of $T_i$. The families are indexed: equal sampled sets retain
their separate indices. The arrays and target are fixed independently
of all four families.

The binomial arguments in \eqref{sep:families} are feasible for
large $n$. To see this, put
\[
 \zeta=\frac12+\alpha\left(k-\frac14\right),\qquad
 Y=\zeta-k=\frac14+(\alpha-1)\left(k-\frac14\right).
\]
Then $e=\zeta g+O(1)$ and $c_i=kg+O(1)$, uniformly in
$k\in[1/6,1/4]$. Since $1<\alpha<2$, $Y$ stays strictly between
$0$ and $1/2$ throughout this interval.

As in Lemma~\ref{lem:entropy}, $h$ denotes binary entropy and
$h_*=h(1/4)$.

\begin{lemma}[Weighted disjointness]\label{lem:separator}
For the input and independent families just specified, there is a
weighted-disjointness check with the following guarantees, for all
sufficiently large $n$.
\begin{enumerate}
\item It returns YES only for a pair satisfying
\eqref{sep:target}. For every fixed pair satisfying that relation,
the probability that the check finds a solution is at least
$1-4e^{-10}$.
\item Preprocessing the families takes at most
$\poly(n)2^{2g}$ operations and words. Thereafter the expected check
work is at most
\[
 (|\mathcal L|+|\mathcal R|)\,u^{22}2^{(1-h_*)m}.
\]
\item Workspace, including preprocessing, is deterministically at most
$O(nC+n^2+2^{2g})$ words. Any execution stopped early has no
larger work or space usage than the complete traversal used in these
bounds.
\end{enumerate}
The probability statements concern the stated ideal uniform families.
Generating them with at most $n^3$ attempts per ordinary rejection
draw also has deterministic cost $\poly(n)2^{2g}$, or reports sampling
failure. It can be coupled to the ideal families with failure
probability $\exp(-\Omega(n^2))$.
\end{lemma}

\begin{proof}
We first check coverage of one fixed pair. We then describe a traversal
whose workspace is bounded for every family outcome. Its time will be
charged to individual occurrences of input records, and an entropy
estimate will bound their expectation.

\paragraph{Coverage of one pair.}
Fix a pair satisfying \eqref{sep:target}. On block $T_i$, a sampled
set must contain its left mask and avoid its right mask. Their union
has size $g/2$, so the separation probability is exactly
\[
 \theta_i=\binom{g/2}{e-c_i}\big/\binom{g}{e}.
\]
Equation~\eqref{sep:families} gives $\nu_i\theta_i\ge10$.
The probability that no family member separates the pair on this
block is at most $(1-\theta_i)^{\nu_i}\le e^{-10}$.
A union bound over the four blocks shows that some tuple of family
indices separates the whole pair with probability at least
$1-4e^{-10}$. This is coverage of a fixed pair; no simultaneous
coverage assertion for every pair is needed.

\paragraph{Preprocessing and the traversal.}
A left record is compatible with index $j$ on block $i$ when
$D\cap T_i\subseteq S_{i,j}$; a right record is compatible when
$(E\cap T_i)\cap S_{i,j}=\varnothing$.
For every mask on each block, store both sorted lists of compatible
indices, one for containment and one for avoidance. Since
$\binom{g/2}{e-c_i}\ge1$,
\[
 \nu_i\le10\cdot2^g+1.
\]
Store a block mask and a family index in one word each. For each block
and each of its $2^g$ masks, the two incidence lists together need at
most $2\nu_i$ index words; their offsets and lengths use only a
constant number of words. Across the four blocks, a flat array of
capacity at most $8\cdot2^g(10\cdot2^g+1)$ suffices for their
contents. Generating the families needs $O(2^g+n)$ words, including
the current permutation. Local-index maps and other control data
use at most $O(n^2)$ words. Thus the families and incidence tables
occupy $O(2^{2g}+n^2)$ words. Filling the lists in increasing index
order costs $\poly(n)2^{2g}$ operations.
A binary search in the stored list decides whether a record has a compatible index in
any given interval of family indices.

Process the blocks in order. For the current block, begin with its
whole index interval $[1,\nu_i]$. Filter both received arrays to
records having a compatible index in this interval. If either filtered
array is empty, finish this branch. Otherwise split a nonsingleton
interval into two nearly equal intervals and visit both children. At
a singleton, fix that family index and enter the next block with the
filtered arrays. After fixing all four indices, sort the two arrays
by actual weight and search for weights summing exactly to $t'$.

Every pair at this last search is disjoint, because its left mask lies
inside the chosen set and its right mask lies outside it on every
block. Thus a reported YES is valid. A fixed pair separated by a tuple
of indices remains in both arrays on the entire path to that tuple.
It therefore reaches an exact-weight search, unless a previous search
has already returned YES. This proves the coverage guarantee.

Use depth-first traversal and retain the parent arrays until their
children finish. Each depth has its own buffers, written only up to
their current lengths. In particular, a small filtered array does
not require clearing a length-$C$ buffer. The combined depth of the
four interval trees is at most
\[
 4+\sum_{i=1}^4\lceil\log\nu_i\rceil\le m+20=O(n).
\]
At each depth at most two arrays of length $C$ are retained, and each
record occupies a constant number of words. The retained-buffer term
is therefore $O(nC)$ words, including the input arrays. Sorting at a
terminal branch uses only $O(C)$ additional words, and stack control
uses $O(n)$ words. Together with the incidence tables this gives the
explicit bound
\begin{equation}\label{sep:workspace}
 O(nC+n^2+2^{2g})\quad\text{writable words}.
\end{equation}
This bound applies to every choice of the families and every feasible
input array. Any other arrays retained by the caller are additional.

\paragraph{Charging the visited occurrences.}
Fix one input occurrence, and let $t_i$ be its number of compatible
indices in family $i$. For a compatible prefix of previously fixed
indices, the occurrence is scanned at the next block's root. If
$t_i=0$, that scan removes it. If $t_i>0$, it can pass the filter only
on paths to its $t_i$ compatible leaves. The union of these paths has
$O(t_i(1+\log\nu_i))$ nodes; scans into their children have the same
order of cost. There are at most $\prod_{r<i}t_r$ compatible
prefixes before block $i$. Its scans at that block are consequently
bounded by a polynomial factor times
\[
 \left(\prod_{r<i}t_r\right)(1+t_i).
\]
The term $1$ is necessary even when $t_i=0$.

An interval predicate costs $O(n)$ operations, including binary
search and local-mask extraction, and each tree has depth $O(n)$.
After the fourth block, sorting and searching cost $O(n)$ operations
per surviving occurrence. These costs give an $O(n^2)$ factor in
the preceding scan count.
Node overhead can be charged to scans because a child or new block is
entered only with nonempty arrays. For sufficiently large $n$, the
complete charge to this input occurrence is at most
\begin{equation}\label{sep:charge}
 u^4\sum_{j=0}^4\prod_{i=1}^j t_i,
\end{equation}
where the empty product equals one. This charge bounds a traversal
which completes all branches even after finding a solution. Stopping
at a YES, an operation limit, or any earlier point can only reduce
its work.

\paragraph{Expected incidences on one block.}
For a fixed left mask of size $c=c_i$, the probability of containment
in a uniform $e$-set is
$\binom{g-c}{e-c}/\binom{g}{e}$. Since
$\nu_i\le11/\theta_i$,
\begin{equation}\label{sep:incidence}
 \E t_i\le
 11\binom{g-c}{e-c}\bigg/\binom{g/2}{e-c}.
\end{equation}
Apply the binomial bounds of Lemma~\ref{lem:entropy}.
Set
\[
 X=1-k,\qquad
 Y=\frac14+(\alpha-1)\left(k-\frac14\right),\qquad
 E(k)=Xh(Y/X)-\frac12h(2Y).
\]
Those bounds show that the logarithm of the
right-hand side of \eqref{sep:incidence} is at most
\begin{equation}\label{sep:entropy-rate}
 gE(k)+O(\log n).
\end{equation}
The bound is uniform in $k\in[1/6,1/3]$. Indeed, $Y$ lies between
$k$ and $1/4$, and $Y/X$ and $2Y$ stay away from zero and one on
this compact interval. Their entropy derivatives are bounded in a
fixed neighborhood. The $O(1)$ rounding of $c$ and $e$ therefore
changes each binomial entropy exponent by only $O(1)$. The lower
binomial bound for the denominator contributes at most
$\log(g/2+1)$ in addition to constants.

For right records use the complements of the sampled sets. Write
\[
 k'=\frac12-k,\qquad c'=g/2-c,\qquad e'=g-e.
\]
The complemented sets are uniform $e'$-sets. Their ideal bias is
$1-\zeta=1/2+\alpha(k'-1/4)$, and rounding again costs only a
bounded amount. The family-size denominator stays the same because
\[
 e'-c'=g/2-(e-c),\qquad
 \binom{g/2}{e'-c'}=\binom{g/2}{e-c}.
\]
Thus \eqref{sep:incidence}--\eqref{sep:entropy-rate} apply to
right records with $k'\in[1/4,1/3]$. It remains to prove one sharp
entropy bound covering both sides.

\paragraph{The exact entropy bound.}
We claim
\begin{equation}\label{sep:sharp}
 E(k)\le1-h_*\qquad(1/6\le k\le1/3).
\end{equation}
Let $F=(\ln2)E$, $Z=X-Y$, and $W=1/2-Y$. Expanding entropy gives
\[
 F=X\ln X-Z\ln Z-\frac12\ln\frac12+W\ln W.
\]
Consequently
\[
 F'=\ln\frac{Z}{X}+(\alpha-1)\ln\frac{Z}{W},
 \qquad
 F''=\frac1X-\frac{\alpha^2}{Z}
                  +\frac{(\alpha-1)^2}{W}.
\]
All denominators are positive even on $[0,1/2]$: explicitly,
\[
 Z=\frac{2+\alpha}{4}-\alpha k,
 \qquad W=\frac\alpha4-(\alpha-1)k,
 \qquad 1<\alpha<2.
\]
The ratio $Z/W$ decreases on this interval, since
\[
 \left(\frac{Z}{W}\right)'=\frac{\alpha-2}{4W^2}<0.
\]
Hence $Z/W\le(2+\alpha)/\alpha$. Also $Z/X\le1$, since
$Y>0$. Multiplying the second derivative by $Z$ now gives
\[
 ZF''=\frac{Z}{X}-\alpha^2+(\alpha-1)^2\frac{Z}{W}
 \le1-\alpha^2+(\alpha-1)^2\frac{2+\alpha}{\alpha}
 =-2+\frac2\alpha<0.
\]
Thus $E$ is strictly concave. At $k=1/4$ we have
$X=3/4$, $Z=1/2$, and $W=1/4$, so
$F'=\ln(2/3)+(\alpha-1)\ln2=0$. Its maximum is therefore
\[
 E(1/4)=\frac34h(1/3)-\frac12=1-h(1/4).
\]
The last identity follows from
$h(1/3)=\log3-2/3$ and $h(1/4)=2-(3/4)\log3$.
This proves \eqref{sep:sharp} without a positive error in its
exponential rate.

Combining \eqref{sep:incidence}--\eqref{sep:sharp}, for either
side and every block,
\[
 \E t_i\le u^4 2^{(1-h_*)g}
\]
for sufficiently large $n$. The four block families are independent,
so expectations of the products in \eqref{sep:charge} factor.
Since $1-h_*>0$ and $m=4g$, each of the five terms in
\eqref{sep:charge} has expectation at most
$u^{20}2^{(1-h_*)m}$. Their sum is at most
$u^{22}2^{(1-h_*)m}$ for large $n$. Summing over input occurrences proves the
expected work bound. The preceding construction already proved the
deterministic space bound.

Finally, a uniform $e$-set is obtained by taking the first $e$
positions of a uniform permutation of its block. Generate each
permutation from fresh successive range draws, with the bounded
implementation of Lemma~\ref{lem:samplers}. Across all four families
there are at most
\[
 4g(10\cdot2^g+1)
\]
such calls. Its conditional per-call bound gives a coupling failure
probability at most $4g(10\cdot2^g+1)2^{-n^3}$, which is
$\exp(-\Omega(n^2))$. The bound uses the ideal product law and
subtracts the coupling failure; it does not condition that law on
successful bounded sampling. Allowing all $n^3$ attempts per draw
adds only a polynomial factor to the generation work, which fits
$\poly(n)2^{2g}$ together with preprocessing.
\end{proof}

\subsection{Reusing the same families}

The check can be called repeatedly with the same four families.
Different calls are then correlated, but their work still adds by
linearity of expectation. The necessary condition is that the full
list of their input arrays be fixed independently of those families.

\begin{corollary}[Reusing the families]\label{cor:separator-reuse}
Fix a finite list of checks satisfying the input conditions of
Lemma~\ref{lem:separator}, with common block parameters, and let
\[
 N_{\mathrm{check}}=\sum_{\text{checks}}(|\mathcal L|+|\mathcal R|).
\]
The list may itself be random, provided that, conditional on the list
and its parameters, the four families have the independent
distributions \eqref{sep:families}. In the ideal experiment, the
expected total check work conditional on the list and its parameters
is at most
\[
 W_0=N_{\mathrm{check}} u^{22}2^{(1-h_*)m}.
\]
Preprocessing is paid once. Processing checks one at a time uses the
check workspace of Lemma~\ref{lem:separator}, in addition to any
storage retained by the caller. For any $0<\delta<1$, total
check work exceeds $W_0/\delta$ with probability at most $\delta$.
Suppose a designated check and a solution pair in it are also chosen,
and that conditional on the list, its parameters, and these choices,
the families still have the independent distributions
\eqref{sep:families}. Then with probability at least
$1-4e^{-10}-\delta$ the pair is covered and the total check work is
at most $W_0/\delta$.
\end{corollary}

\begin{proof}
An empty list has $N_{\mathrm{check}}=0$ and zero check work. Otherwise
condition on the complete list and its parameters, and apply Lemma~\ref{lem:separator}
to each input occurrence and sum expectations. Markov's inequality
gives the work bound. For the last assertion, condition additionally
on the designated check and pair. The family law is unchanged, so
the same expected-work bound holds and the pair has coverage
probability at least $1-4e^{-10}$. A union bound intersects the
coverage and work events; those events need not be independent.
\end{proof}

In the Subset Sum application, continue parent generation as if checks
never returned YES and as if no operation cap were present, retaining
all dictionary overflow rules. This defines the complete check list
without consulting the separating families. Conditional on the layout,
their sampled product law is independent of all parent-generation
choices. It therefore remains unchanged after fixing that list and a
designated check and witness selected from it using only those choices.
Corollary~\ref{cor:separator-reuse} applies once
Section~\ref{sec:resources} bounds the total input length of this list.
The actual run visits only an initial portion of its work; the list
itself is never stored, and coverage is required only for the selected
pair.
